\section*{\texorpdfstring{\S\CurrentExerciseGroup}{第{\CurrentExerciseGroup}节}}
\phantomsection
\addcontentsline{toc}{section}{第{\CurrentExerciseGroup}节习题}

\begin{enumerate}
\def\labelenumi{\arabic{enumi})}
\item
  设 T 是 R 的紧区间 \([0,1]\) ，\(\mu\) 是 T 上的 Lebesgue 测度。设 X 是把区间 \([0,1]\) 赋以离散拓扑所得的局部紧空间。映射 \(t \mapsto \lambda_{t} = \varepsilon_{t}\) （T 到 \(\mathcal{M}(X)\) 中的映射）是标量本质可积的，其积分为 \(\nu = 0\) 。对 X 上等于 1 的常值函数 f，命题 3 的公式 (6) 不成立。由此推出映射 \(t \mapsto \lambda_{t}\) 不是 \(\mu\) -适宜的。
\item
  设 T 和 \(\mu\) 是 Ch. IV, §1 习题 5 中定义的局部紧空间和测度 \(\geqslant 0\) 。取 X 为集 T 赋以由 \(R^{2}\) 诱导的拓扑所得的空间，于是 X 是局部紧的且在无穷远处可数。设 \(t \mapsto \lambda_{t}\) 是 T 到 \(\mathcal{M}(X)\) 中的映射，满足 \(\lambda_{t} = \varepsilon_{t}\) （当 \(t = (0, y)\) ），且 \(\lambda_{t} = \varepsilon_{t}/n^{3}\) （当 \(t = (1/n, k/n^{2})\) ）。证明映射 \(t \mapsto \lambda_{t}\) 是 \(\mu\) -适宜的、淡 \(\mu\) -可测的、但不淡连续，且若 \(\nu = \int \lambda_{t} d\mu(t)\) 而 \(f = \varphi_{I}\) ，其中 I 是 \(R^{2}\) 中以 0 为中心、边长为 2 的正方形，则
\end{enumerate}

\[
\int^ {*} f (x) d \nu (x) <   \int^ {*} d \mu (t) \int^ {\bullet} f (x) d \lambda_ {t} (x) = + \infty .
\]

\begin{enumerate}
\def\labelenumi{\arabic{enumi})}
\setcounter{enumi}{2}
\tightlist
\item
  设 \(\mathrm{T}, \mathrm{T}'\) 是两个局部紧空间，\(\mu\) 是 \(\mathrm{T}\) 上的一个正测度，\(\mu'\) 是 \(\mathrm{T}'\) 上的一个正测度。在乘积空间 \(\mathrm{X} = \mathrm{T} \times \mathrm{T}'\) 上，置 \(\lambda_t = \varepsilon_t \otimes \mu'\) 、\(\lambda_{t'}' = \mu \otimes \varepsilon_{t'}\) （对 \(t \in \mathrm{T}\) 、\(t' \in \mathrm{T}'\) ）。映射 \(t \mapsto \lambda_t\) （相应地，\(t' \mapsto \lambda_{t'}'\) ）——\(\mathrm{T}\) （相应地，\(\mathrm{T}'\) ）到 \(\mathcal{M}(\mathrm{X})\) 中的映射——是淡连续且 \(\mu\) -适宜的（相应地，\(\mu'\) -适宜的），且
\end{enumerate}

\[
\nu = \mu \otimes \mu^ {\prime} = \int \lambda_ {t} d \mu (t) = \int \lambda_ {t ^ {\prime}} ^ {\prime} d \mu^ {\prime} (t ^ {\prime})
\]

（Ch. III, §4, No.~1；参见 §8）。取 T 为 R 的区间 {[}0,1{]}，取 \(\mu\) 为 Lebesgue 测度，取 T′ 为赋以离散拓扑的区间 {[}0,1{]}，并取 \(\mu'\) 为 T′ 上由每点质量 +1 定义的测度。设 \(f\) 是 X 的“对角线”（由点 \((t,t)\) （其中 \(t\) 取遍 {[}0,1{]}）组成之集）的特征函数。证明 \(f\) 是上半连续的，且

\[
1 = \int^ {*} d \mu (t) \int^ {*} f (x) d \lambda_ {t} (x) <   \int^ {*} f (x) d \nu (x) = + \infty
\]

以及

\[
0 = \int^ {\bullet} f (x) d \nu (x) <   \int^ {\bullet} d \mu (t) \int^ {\bullet} f (x) d \lambda_ {t} (x) = 1.
\]

¶ 4) 设空间 T、T′、X 和测度 \(\mu, \mu', \nu\) 的意义与习题 3 中相同，考虑乘积空间 Y = T × X = T × (T × T′)，以及 Y 上的乘积测度 \(\varpi = \mu \otimes \nu\) 。若对 \(t \in \mathrm{T}\) 置 \(\rho_t = \varepsilon_t \otimes \nu\) ，则有 \(\varpi = \int \rho_t d\mu(t)\) 。设 H 是 T 的一个对 \(\mu\) 不可测的子集（Ch. IV, §4, 习题 8），并设 A 是由点 \((t_1, t_2, t_3)\) ——满足条件 \(t_1 = t_3\) 、\(t_2 \in \mathrm{H}\) ——组成的 Y 的子集。证明特征函数 \(\varphi_A\) 对 \(\varpi\) 是局部可忽略的，但 \(\varphi_A\) 不是 \(\rho_t\) -可测的（对任何 \(t \in \mathrm{T}\) 的值）。

\begin{enumerate}
\def\labelenumi{\arabic{enumi})}
\setcounter{enumi}{4}
\item
  给出一个 \(\mu\) -适宜族 \(t \mapsto \lambda_t\) 和一个数值函数 \(f\) （定义在 \(X\) 上）的例子，使得 \(f\) 是 \(\lambda_t\) -可测的（对所有 \(t \in T\) ），但对测度 \(\nu = \int \lambda_t d\mu(t)\) 不可测（取 \(X = T\) 和 \(\lambda_t = \varepsilon_t\) （对所有 \(t \in T\) ））。
\item
  \begin{enumerate}
  \def\labelenumii{\alph{enumii})}
  \tightlist
  \item
    证明：在 §3 中涉及淡连续映射 \(t \mapsto \lambda_t\) 这一概念的所有结果中，都可以把这一假设换成下列假设：对每个函数 \(g \in \mathcal{K}_+(X)\) ，函数 \(t \mapsto \lambda_t(g)\) 在 \(T\) 上是下半连续的。
  \end{enumerate}
\end{enumerate}

\begin{enumerate}
\def\labelenumi{\alph{enumi})}
\setcounter{enumi}{1}
\tightlist
\item
  考虑一个标量本质可积映射 \(t \mapsto \lambda_t\) （\(T\) 到 \(\mathcal{M}_+(X)\) 中的映射），它满足下列条件：
\end{enumerate}

对 T 的每个紧子集 K 和每个 \(\varepsilon > 0\) ，存在紧集 \(K_{1} \subset K\) ，使得 \(\mu(K - K_{1}) \leqslant \varepsilon\) ，且在 \(K_{1}\) 上，所有函数 \(t \mapsto \langle f, \lambda_{t} \rangle\) （其中 f 取遍 \(\mathcal{K}_{+}(X)\) ）的限制都是下半连续的。

证明映射 \(t \mapsto \lambda_t\) 是 \(\mu\) -适宜的。

\begin{enumerate}
\def\labelenumi{\arabic{enumi})}
\setcounter{enumi}{6}
\tightlist
\item
  \begin{enumerate}
  \def\labelenumii{\alph{enumii})}
  \tightlist
  \item
    设 \(\Lambda : t \mapsto \lambda_t\) 是 T 到 \(\mathcal{M}_{+}(X)\) 中的一个标量本质可积映射，且 \(\nu = \int \lambda_t d\mu(t)\) 。证明对每个下半连续函数 \(f \geqslant 0\) （定义在 \(X\) 上），
  \end{enumerate}
\end{enumerate}

\[
\int^ {\bullet} f (x) d \nu (x) \leqslant \int^ {\bullet} d \mu (t) \int^ {\bullet} f (x) d \lambda_ {t} (x).
\]

\begin{enumerate}
\def\labelenumi{\alph{enumi})}
\setcounter{enumi}{1}
\tightlist
\item
  设 \(\Lambda\) 是 \(\mu\) -预适宜的，用 \(\mu'\) 表示一个正测度 \(\leqslant \mu\) ，并写 \(\mu = \mu' + \mu''\) 、\(\nu' = \int \lambda_t d\mu'(t)\) 。由 \(a\) ) 推出，对每个下半连续函数 \(f \geqslant 0\) （\(\nu\) -可积的），
\end{enumerate}

\[
\int^ {\bullet} f (x) d \nu^ {\prime} (x) = \int^ {\bullet} d \mu^ {\prime} (t) \int^ {\bullet} f (x) d \lambda_ {t} (x).
\]

把这一结果推广到下半连续函数 \(f \geqslant 0\) （\(\nu\) -适度的）。由此推出，若 \(\nu\) 是一个适度测度（特别地，若 X 在无穷远处可数），则 \(\Lambda\) 是 \(\mu\) -适宜的。

\begin{enumerate}
\def\labelenumi{\arabic{enumi})}
\setcounter{enumi}{7}
\tightlist
\item
  设 \(\Lambda : t \mapsto \lambda_t\) 是 \(T\) 到 \(\mathcal{M}_+(X)\) 中的一个映射。
\end{enumerate}

\begin{enumerate}
\def\labelenumi{\alph{enumi})}
\item
  设 \(\mu_{1},\ldots,\mu_{n}\) 是 T 上的测度，且 \(\mu=\mu_{1}+\cdots+\mu_{n}\) 。为使 \(\Lambda\) 是 \(\mu\) -适宜的，必须且只须 \(\Lambda\) 对每个 i 都是 \(\mu_{i}\) -适宜的（利用分解引理）。
\item
  设 \(\mu\) 是递增有向族 \((\mu_i)_{i \in I}\) 的上确界。为使 \(\Lambda\) 是 \(\mu\) -适宜的，必须且只须 \(\Lambda\) 是 \(\mu_i\) -适宜的（对所有 \(i \in I\) ）且是标量本质 \(\mu\) -可积的。
\item
  设 \(\mu\) 是正测度组成的可和族 \((\mu_j)_{j \in \mathbb{J}}\) 之和。为使 \(\Lambda\) 是 \(\mu\) -适宜的，必须且只须 \(\Lambda\) 是 \(\mu_j\) -适宜的（对所有 \(j \in \mathbb{J}\) ），且 \(\Lambda\) 是标量本质 \(\mu\) -可积的。
\item
  为使 \(\Lambda\) 是 \(\mu\) -适宜的，必须且只须 \(\Lambda\) 是标量本质 \(\mu\) -可积的，且 \(\Lambda\) 是 \(\mu'\) -预适宜的（对每个满足 \(\mu' \leqslant \mu\) 的具紧支集测度）。
\end{enumerate}

\begin{enumerate}
\def\labelenumi{\arabic{enumi})}
\setcounter{enumi}{8}
\tightlist
\item
  \begin{enumerate}
  \def\labelenumii{\alph{enumii})}
  \tightlist
  \item
    设 \(\mathrm{T}\) 和 \(\mathrm{X}\) 是两个局部紧空间，\(\mathcal{C}_{+}(\mathrm{T})\) 是 \(\mathrm{T}\) 上所有正连续函数组成的凸锥，而 \(\mathrm{V}\) 是 \(\mathcal{H}_{+}(\mathrm{T})\) 到 \(\mathcal{C}_{+}(\mathrm{T})\) 中的一个映射，具有下列性质：
  \end{enumerate}
\end{enumerate}

\[
\begin{array}{r l} \mathrm{V} (f + g) = \mathrm{V} f + \mathrm{V} g & \text {if} f, g \in \mathcal {K} _ {+} (\mathrm{X}) \\ \mathrm{V} (t f) = t (\mathrm{V} f) & \text {if} f \in \mathcal {K} _ {+} (\mathrm{X}), t \in \mathbf {R} _ {+}. \end{array}
\]

证明存在唯一的扩散 \(\Lambda\) （\(T\) 到 \(X\) 中的扩散），使得 \(\Lambda f = Vf\) （对所有 \(f \in \mathcal{K}_{+}(X)\) ）。证明：若把锥 \(\mathcal{C}_{+}(T)\) 换成 \(T\) 上正的且局部有界的下半连续函数组成的锥，这一结果仍成立（参见习题 6）。

\begin{enumerate}
\def\labelenumi{\alph{enumi})}
\setcounter{enumi}{1}
\tightlist
\item
  设 X 的拓扑有一个可数基。证明：若把锥 \(\mathcal{C}_{+}(\mathrm{T})\) 换成 T 上正的且局部有界的普遍可测函数组成的锥，上述结果仍成立。
\end{enumerate}

\begin{enumerate}
\def\labelenumi{\arabic{enumi})}
\setcounter{enumi}{9}
\item
  设 \(\mathrm{T}, \mathrm{X}, \mathrm{Y}\) 是三个在无穷远处可数的局部紧空间，\(\Lambda : t \mapsto \lambda_t\) 是 \(\mathrm{T}\) 到 \(\mathrm{X}\) 中的一个扩散，\(\mathrm{H}: x \mapsto \eta_x\) 是 \(\mathrm{X}\) 到 \(\mathrm{Y}\) 中的一个扩散；称 \(\Lambda\) 与 \(\mathrm{H}\) 是可复合的，如果函数 \(\Lambda(\mathrm{Hg})\) 在 \(\mathrm{T}\) 上局部有界（对每个 \(g \in \mathcal{K}_+(\mathrm{Y})\) ）。证明：若 \(\Lambda\) 与 \(\mathrm{H}\) 可复合，则 \(\lambda_t\) 属于 \(\mathrm{H}\) 的定义域（对所有 \(t \in \mathrm{T}\) ），且映射 \(t \mapsto \lambda_t\mathrm{H}\) 是 \(\mathrm{T}\) 到 \(\mathrm{Y}\) 中的一个扩散。把公式 (15)（命题 13）推广到这一情形。
\item
  设 \(t \mapsto \lambda_t\) 是一个 \(\mu\) -适宜映射（\(T\) 到 \(\mathcal{M}_+(X)\) 中的映射），满足下列条件：对每个紧子集 \(K\) （\(T\) 的子集），存在紧子集 \(L_K\) （\(X\) 的子集），使得 \(\operatorname{Supp}(\lambda_t) \subset L_K\) （对所有 \(t \in K\) ）。
\end{enumerate}

\begin{enumerate}
\def\labelenumi{\alph{enumi})}
\tightlist
\item
  证明：对每个函数 \(f \geqslant 0\) （定义在 \(X\) 上），
\end{enumerate}

\[
\int^ {\bullet} f (x) d \nu (x) \geqslant \int^ {\bullet} d \mu (t) \int^ {*} f (x) d \lambda_ {t} (x)
\]

（利用命题 3 a)）。

\begin{enumerate}
\def\labelenumi{\alph{enumi})}
\setcounter{enumi}{1}
\item
  若 \(f \geqslant 0\) 是局部 \(\nu\) -可忽略的，则由这样的 \(t \in \mathbb{T}\) 组成之集——使 \(f\) 不是局部 \(\lambda_t\) -可忽略的——是局部 \(\mu\) -可忽略的。
\item
  若 \(f\) 是一个 \(\nu\) -可测映射（\(X\) 到一个拓扑空间 \(X\) 中的映射），则由这样的 \(t \in T\) 组成之集——使 \(f\) 不是 \(\lambda_t\) -可测的——是局部 \(\mu\) -可忽略的。
\item
  对每个 \(\nu\) -可测函数 \(f \geqslant 0\) ，映射 \(t \mapsto \int^{*} f d\lambda_{t}\) 是 \(\mu\) -可测的，且
\end{enumerate}

\[
\int^ {\bullet} f (x) d \nu (x) = \int^ {\bullet} d \mu (t) \int^ {*} f (x) d \lambda_ {t} (x).
\]

\begin{enumerate}
\def\labelenumi{\alph{enumi})}
\setcounter{enumi}{4}
\tightlist
\item
  设 \(\mathbf{f}\) 是定义在 \(X\) 上、取值于一个 Banach 空间或 \(\overline{\mathbf{R}}\) 中的本质 \(\nu\) -可积函数。证明：由这样的 \(t \in T\) 组成之集——使 \(\mathbf{f}\) 不是 \(\lambda_t\) -可积的——是局部 \(\mu\) -可忽略的，函数 \(t \mapsto \int \mathbf{f}(x) d\lambda_t(x)\) （对 \(\mu\) 局部几乎处处有定义）是本质 \(\mu\) -可积的，且
\end{enumerate}

\[
\int \mathbf {f} (x) d \nu (x) = \int d \mu (t) \int \mathbf {f} (x) d \lambda_ {t} (x).
\]

\setcounter{exercisegroup}{4}
\edef\CurrentExerciseGroup{\arabic{exercisegroup}}
